%% kapa: exam=1 round=30 year=2019 period=1 form=A subject=economics subject_ko=경제학원론 seq=30 number=70
%% source: Q-Net 감정평가사 기출문제 (cst00309) artlSeq=5205205 file=extracted/1008/5205205/2019년 제30회 감정평가사 1차시험 문제지/2019년 제30회 감정평가사 1차시험 문제지/2019년도 제30회 감정평가사 1차 1교시 A형.hwp
%% conversion: kordoc 4.24.51 -> tools/gen1.py (auto); review=proofread-sonnet; flags=none
\begin{question}{70}
리카도 대등정리(Ricardian equivalence theorem)는 정부지출의 재원조달 방식에 나타나는 변화가 민간부문의 경제활동에 아무런 영향을 주지 못한다는 것이다. 이 정리가 성립하기 위한 가정으로 옳은 것을 모두 고른 것은?
\begin{framedtext}
ㄱ. 유동성 제약 \par
ㄴ. 경제활동인구 증가율 양(+)의 값 \par
ㄷ. 일정한 정부지출수준과 균형재정 \par
ㄹ. ‘합리적 기대’에 따라 합리적으로 행동하는 경제주체
\end{framedtext}
\begin{choices}
  \item ㄱ, ㄴ
  \item ㄴ, ㄷ
  \item ㄷ, ㄹ
  \item ㄱ, ㄷ, ㄹ
  \item ㄴ, ㄷ, ㄹ
\end{choices}
\end{question}
\answer{3}
